% =====================================================================
% Seção 1 — Introdução
% =====================================================================
\section{Introdução}\label{sec:intro}

A fronteira entre o que a inteligência artificial produz \emph{sozinha} e o que
ela obtém \emph{do mundo} tornou-se, nos últimos três anos, uma questão central
tanto da engenharia quanto da teoria. Dois fenômenos simultâneos a tornaram
urgente. De um lado, a escala dos \emph{large language models} (LLMs) apróxima-se
do limite de dados humanos disponíveis: as \emph{scaling laws} empíricas de
Kaplan \textit{et al.}~\parencite{Kaplan2020Scaling} e a revisão
Chinchilla~\parencite{Hoffmann2022Chinchilla} estabeleceram o dado, ao lado de
parâmetros e computação, como um dos três insumos físicos que determinam a perda
de um modelo, e o regime de restrição a dados
(\emph{data-constrained})~\parencite{Muennighoff2023DataConstrained} tornou-se o
cenário dominante de planejamento. De outro, a prática industrial passou a
empregar \emph{dados sintéticos} --- respostas, raciocínios e problemas gerados por
modelos --- como combustível central do pós-treino.

Nesse encontro nasce a pergunta que este artigo investiga: \textbf{existe um
fator limítrofe entre produzir IA com IA?} A questão não é meramente filosófica.
Ela tem consequências operacionais diretas: se um modelo pode, em princípio,
gerar o seu próprio sinal de melhoria de forma ilimitada, então o caminho para a
auto-melhoria recursiva seria uma reta sem teto; se, ao contrário, há um limite
informacional e de confiabilidade, então todo loop de auto-produção precisa de
uma âncora externa para não degenerar.

Evidências recentes e de alto impacto apontam para a segunda leitura.
Shumailov \textit{et al.}~\parencite{Shumailov2024Collapse}, em \emph{Nature},
demonstraram que modelos treinados recursivamente sobre dados gerados por
modelos \emph{colapsam}: a distribuição aprendida perde os modos de baixa
probabilidade, a variância colapsa e o modelo converge para uma singularidade
estatística. Borji~\parencite{Borji2024CollapseNote} mostrou que o fenômeno é um
efeito estatístico fundamental, provavelmente inevitável em loop fechado. Em
paralelo, a literatura de auto-melhoria documenta um gargalo de confiabilidade:
\emph{self-play} instável~\parencite{SelfPlayGuidance2026} e um
\emph{discovery--reliability gap} em agentes de pesquisa
recursiva~\parencite{RSIBench2026}.

Este artigo contribui ao integrar essas três linhas --- colapso, \emph{scaling
laws} e auto-melhoria --- com um \textbf{estudo de caso empírico} conduzido no
\emph{OpenCode Ecosystem Core}: um ecossistema multiagente metacognitivo no qual
ciclos de pesquisa (R458--R460) avaliaram a diversificação pós-ranqueamento em
\emph{Retrieval-Augmented Generation} (RAG). O estudo de caso é instrutivo porque
encarna, em microcosmo e com dados primários auditáveis, exatamente o \emph{trade-off
} limítrofe que a teoria prevê: ganhar na fronteira de uma métrica (diversidade)
custa fidelidade a outra (relevância).

A resposta que defendemos, ancorada nas evidências, é equilibrada e honesta:
\textbf{existe um fator limítrofe estrutural, mas não absoluto}. Ele separa a
produção de IA em \emph{loop fechado auto-suficiente} (que colapsa) da produção
\emph{com âncora externa de dados reais} (que progride). O restante do artigo
estrutura-se assim: a Seção~\ref{sec:fund} estabelece a fundamentação teórica; a
Seção~\ref{sec:metodo} conceitua formalmente o fator limítrofe e descreve o
método do estudo de caso; a Seção~\ref{sec:caso} apresenta os resultados
empíricos; a Seção~\ref{sec:disc} discute a síntese; e a Seção~\ref{sec:conc}
responde à pergunta central, declara os limites e evita alegações excessivas.
